\section{Introduction}
\label{sec:introduction}

The automated detection and morphological classification of lunar surface features is a critical capability for modern planetary science. High-resolution orbital imagery from the Lunar Reconnaissance Orbiter (LRO) now covers the Moon at resolutions down to half a metre per pixel, producing volumes of data that far exceed manual analysis capacity. Impact craters, wrinkle ridges, lobate scarps, lava-tube skylights, and other geomorphological features encode the thermal, tectonic, and impact history of the lunar surface, and their reliable identification is a prerequisite for geological mapping, mission planning, and the selection of future landing sites.

The computer-vision literature offers several distinct formulations for this task, and it is not obvious a priori which is best suited to planetary terrain. Semantic segmentation assigns a label to every pixel and naturally represents extended, coverage-like structures; instance segmentation additionally separates individual objects, which matters when features must be counted; single-stage object detection trades pixel precision for speed by predicting bounding boxes directly; and classical image processing (thresholding, morphology, ridge filters) remains attractive for its interpretability and negligible computational cost. Published comparisons between these families are typically confounded by differences in datasets, train/test splits, metrics, and operating thresholds, making their conclusions difficult to transfer.

This work removes those confounds by comparing one representative of each family---a thresholding-and-morphology baseline, a compact U-Net (SmallUNet), an adapted Mask R-CNN, and YOLOv8---on the \emph{same} data, the \emph{same} leakage-free spatial split, and the \emph{same} metrics. The comparison is organised per feature class rather than per architecture, because the dataset itself makes a single aggregate ranking meaningless: its label channels range from an impact-crater layer covering 82\% of all pixels to classes with fewer than one positive pixel in $10^{4}$. Two methodological experiments frame the comparison. First, we show that the saturated crater channel reduces naive scores to a base-rate artefact---a predict-everything baseline attains F1 $= 0.90$---and adopt prevalence-robust metrics accordingly. Second, we quantify the train--validation leakage induced by 50\%-overlapping tiles: identical detectors are retrained under a random tile split and under a spatial block split, isolating how much of a reported score is memorisation rather than generalisation.

Each method section documents the architecture, its training regime, and the design decisions specific to lunar imagery (Sects.~\ref{sec:classical}--\ref{sec:yolo}); the shared-protocol comparison and the split-protocol study follow (Sect.~\ref{sec:comparison}); the pipelines are then deployed on the lunar south pole, a region without ground truth, as a transfer diagnostic (Sect.~\ref{sec:southpole}). Beyond the specific rankings, the contribution of this study is methodological: on planetary datasets with extreme prevalence structure, the evaluation protocol---split geometry, metric choice, threshold calibration---moves reported numbers by more than any architectural decision we tested, and must therefore be fixed before architectures are compared.
